%File: paperA_draft_v2.tex
% Paper A -- draft v3 (restructured 2026-08-11 against five ML-paper-writing guides:
%   Nanda "Highly Opinionated Advice on How to Write ML Papers";
%   ICML 2022 Best Practices; Foerster "How to ML Paper"; Perez "Easy Paper Writing
%   Tips"; Nanda/Perez "Tips for Empirical Alignment Research").
% v2 -> v3 changes are recorded in paper/paperA_v3_RESTRUCTURE.md.
%
% Basis: inputs/Gleb_overleaf_paper.txt (Iteration B, frozen 2026-01-26)
%        restructured per graph_analysis/phase2_results/paperA_positive_outline_2026-06-26.md
% Every number carries a source comment: SRC = on-disk receipt in graph_analysis/phase2_results/,
% and is regression-checked by graph_analysis/experiment_paper_claim_audit.py.
% Markers: [GAP: ...] missing content, [CITE: ...] bib entry needed.
%
% CLASS NOTE: deliberately NOT a conference class. Conference styles (aaai, acmart,
% IEEEtran) carry publisher branding and licence conditions that make them a poor fit
% for an arXiv preprint that has not been accepted anywhere. This is stock `article`
% in two-column mode, which produces workshop-style formatting with no such strings
% attached. Switching to a venue class later is a preamble-only change.
\documentclass[10pt,twocolumn,letterpaper]{article}

\usepackage[T1]{fontenc}
\usepackage[utf8]{inputenc}
\usepackage{times}
\usepackage[margin=0.75in,columnsep=0.25in]{geometry}
\usepackage{graphicx}
\usepackage{booktabs}
\usepackage{amsmath,amssymb}
\usepackage[hyphens]{url}
\usepackage{natbib}
\usepackage{caption}
\usepackage{abstract}
\usepackage{titlesec}
\usepackage{authblk}
% hyperref + cleveref per Foerster's formatting checklist. hyperref must load late;
% cleveref must load after hyperref.
\usepackage[hidelinks,backref=page]{hyperref}
\usepackage[capitalise,noabbrev]{cleveref}

% compact workshop-style headings
\titleformat{\section}{\normalfont\large\bfseries}{\thesection}{1em}{}
\titleformat{\subsection}{\normalfont\normalsize\bfseries}{\thesubsection}{1em}{}
\titlespacing*{\section}{0pt}{1.4ex plus .2ex}{0.8ex plus .1ex}
\titlespacing*{\subsection}{0pt}{1.2ex plus .2ex}{0.6ex plus .1ex}
\setlength{\parskip}{0pt}
\setlength{\bibsep}{0pt plus 0.3ex}
\renewcommand{\abstractnamefont}{\normalfont\bfseries}
\renewcommand{\abstracttextfont}{\normalfont\small}
\frenchspacing
\pagestyle{plain}

\title{\vspace{-2em}\textbf{Reasoning Chains, Not Knowledge Graphs:\\
Mechanism-Level Extraction from the AI Safety Literature\\
and the Limits of Graph-Analytic Inference}\vspace{-0.5em}}

% [GAP: author list + affiliations. Contributors: Martin Leitgab, Gleb Maksimov,
%  Sai (validation), Mike (judge data), Axel (related work), Jeffrey Parks (review).
%  Confirm order + affiliations before submission. Using authblk:
%    \author[1]{Name}  \affil[1]{Institution, \texttt{email}}]
\author{Anonymous Submission}
\date{}

\begin{document}

% Title and abstract span both columns (the standard twocolumn-article idiom).
% Without this wrapper they would be trapped in the left column.
\twocolumn[
  \begin{@twocolumnfalse}
  \maketitle
  \begin{abstract}
  \noindent
Directories of AI safety research catalogue papers by metadata; embedding maps
place them by topical proximity. Neither records \emph{by what mechanism} a
proposed intervention is argued to reduce a specific risk, because that argument
is internal to one document rather than a relation between documents. We prompt a
reasoning LLM to read 11{,}779 documents from the Alignment Research Dataset and
emit, per paper, the chain from a named risk through problem analysis,
theoretical insight, design rationale, implementation mechanism and validation
evidence to a proposed intervention, giving 200{,}525 typed nodes. We defend
three claims. The chains carry argument structure rather than a filled template:
87.4\% of 2{,}772 quality-controlled de-duplicated chains contain all five stages
while length spans 4--16 nodes, and an LLM-as-judge study finds omission, not
fabrication, to be the dominant error. Chains are a usable unit of analysis: a
risk--intervention query returns one stored chain with its citation, though only
one document in six yields a complete chain and only 2.5\% of extracted
interventions are described as deployed. But the aggregate graph built from the
same extraction does not support the network analysis it invites: merging
near-duplicate nodes manufactures a centrality hub at 90$\times$ the next node
that excluding similarity edges fails to repair, and two previously reported
findings on this corpus fail to reproduce. We release the pipeline, the chain
dataset and the audit scripts that regenerate every number here.
% All abstract figures audited against raw data 2026-08-11 --
% experiment_paper_claim_audit.json, 42/42 numeric claims PASS (2026-08-11, v3).
% The unreproducible 88% figure is REMOVED from the abstract (2026-08-11).
% See experiment_race_top100_rederive_report.json and paperA_v2_AUDIT.md finding P5.
  \end{abstract}
  \vspace{0.6em}
  \end{@twocolumnfalse}
]

% Figure 1 spans both columns and lands on page 1, per Foerster and Perez.
\begin{figure*}[t]
\centering
\includegraphics[width=\textwidth]{figure1_dataset.png}
\caption{What the extraction produced, and what it does not contain.
\textbf{(A)} Length distribution of the 2{,}772 de-duplicated quality-controlled
chains. The canonical risk$\rightarrow$intervention chain has seven nodes
(orange); the spread from 4 to 16 nodes is the evidence that extraction follows
each paper's argument instead of filling a fixed template, since 25.8\% of chains
carry extra nodes for some stage and 10.2\% have genuine gaps.
\textbf{(B)} LLM-assessed maturity of all 36{,}959 extracted interventions:
roughly six in seven sit at conceptual or early-stage maturity and 2.5\% are
described as deployed --- a field-state descriptive that paper metadata cannot
yield.
\textbf{(C)} Corpus yield: 1{,}868 of the 11{,}779 documents (15.9\%) produce any
complete high-confidence chain, and those documents contribute every chain in
panel A. Panel B is an LLM-assigned attribute that our judge study does not
score; we read it as coarse composition rather than as a measured rate
(\cref{sec:claim2}).}
\label{fig:dataset}
\end{figure*}
% SRC: graph_analysis/experiment_figure1.py, which plots only from committed receipts
%      (experiment_paper_claim_audit.json, experiment_dataset_strength_report.json,
%      experiment_dataset_strength_report.json). No PKL checkpoint required.

%%%%%%%%%%%%%%%%%%%%%%%%%%%%%%%%%%%%%%%%%%%%%%%%%%%%%%%%%%%%%%%%%%%%%%%%
\section{Introduction}

The AI safety literature grows faster than any individual can track, and the
question a practitioner asks of it is mechanistic: \emph{if I care about risk
$R$, what has been proposed to reduce it, by what mechanism, and how far has that
proposal been validated?} Existing resources cannot answer this. Curated
directories such as the Alignment Research Dataset (ARD) and Stampy index papers
by metadata. Embedding-based field maps place papers near topically similar
papers. In both cases the unit of organisation is the document, so the argument
inside the document --- the chain from a risk to an intervention --- is not
represented and cannot be queried, aggregated, or compared across papers.

Recovering that argument is harder than the usual scientific-information-extraction
task. The relation of interest connects concepts that the same paper introduces,
it spans several reasoning steps rather than one predicate, and no annotated
corpus of such chains exists to train or evaluate against.

We prompt a reasoning LLM (\texttt{o3}) to read each paper and emit a typed
graph: \textsc{risk} and \textsc{intervention} endpoints connected through five
intermediate concept subtypes --- \emph{problem analysis}, \emph{theoretical
insight}, \emph{design rationale}, \emph{implementation mechanism},
\emph{validation evidence} --- with per-edge confidence and per-intervention
maturity and lifecycle attributes. Across the ARD corpus this produces 200{,}525
nodes over 11{,}779 documents.

\paragraph{Thesis.} For a corpus of this kind, \emph{the per-paper reasoning
chain is a reliable unit of analysis and the aggregate graph is not}. We defend
that sentence with three claims, each with its own evidence section.

\paragraph{Claims and evidence.}
\begin{enumerate}
\setlength{\itemsep}{2pt}
\item[\textbf{C1}] \textbf{The extracted chains carry argument structure rather
than a filled-in schema.} 87.4\% of 2{,}772 quality-controlled de-duplicated
chains contain all five intermediate reasoning stages while length spans 4--16
nodes (\cref{fig:dataset}A); an LLM-as-judge study over 100 papers with three
independent meta-graders finds omission, not fabrication, to be the dominant
failure (\cref{sec:claim1}).
\item[\textbf{C2}] \textbf{Chains are a usable unit of analysis, and reading the
corpus at chain level exposes what it holds.} A risk--intervention query returns
one stored chain with its source citation and no clustering step; the same view
shows 2.5\% of interventions described as deployed and a corpus yield of one
complete chain per six documents (\cref{fig:dataset}B,C; \cref{sec:claim2}).
\item[\textbf{C3}] \textbf{The aggregate graph does not support the network
analysis it invites.} Merging near-duplicate nodes manufactures a centrality hub
at 90$\times$ the next node, and two findings previously reported on this corpus
fail to reproduce under re-derivation (\cref{sec:claim3}).
\end{enumerate}

C3 is a negative result, and we report it as a first-class contribution rather
than as a caveat. The quantities that ecosystem-level claims about this
literature are usually built from --- centrality rankings, per-cluster
risk-to-intervention ratios, co-occurrence measured on a centrality-selected head
--- are dominated on this corpus by construction artefacts. Consequently we make
no ecosystem-level claim ourselves.

%%%%%%%%%%%%%%%%%%%%%%%%%%%%%%%%%%%%%%%%%%%%%%%%%%%%%%%%%%%%%%%%%%%%%%%%
\section{Related Work}
% [GAP: Axel started this section. The five strands below are the agreed scope and
%  all citations resolve; a final read-through for flow is the remaining work.]

\paragraph{Knowledge graphs from scientific literature.}
Prior work extracts entities and relations from scientific text, typically
targeting citation, method--task--dataset, or claim structures.
\citet{luan2018scierc} is representative, jointly identifying entities, relations
and coreference to build a scientific knowledge graph, and the task has a long
methodological history \citep{bach2007relation}. Our schema differs in that the
relation of interest is an \emph{argument} internal to one paper --- risk to
intervention through explicit reasoning stages --- rather than a bibliographic or
method-level relation.

\paragraph{Agentic and self-organising knowledge graphs.}
A parallel line of work has models build and extend knowledge graphs
autonomously. \citet{buehler2025agentic} grows a scale-free knowledge network
through iterative model feedback, and \citet{rasmussen2025zep} uses a temporal
knowledge graph as agent memory. Those systems optimise a graph for an agent's
downstream reasoning; we optimise for fidelity to what a specific paper argued,
which is why our edges are confidence-scored assertions traceable to one source
document rather than model-generated links.

\paragraph{AI safety literature resources.}
% [CITE: AI Safety Atlas -- no canonical reference resolved; cite URL inline if used.]
The Alignment Research Dataset \citep{kirchner2022ard}, maintained by the Stampy /
AI Safety Info project, and hand-curated topic maps such as the AI Safety Atlas
organise the field at the level of documents or topics. The MIT AI Risk Repository
\citep{slattery2026airisk} provides a curated risk vocabulary --- 1{,}725 risks
distilled from 74 frameworks into causal and domain taxonomies. These complement
rather than compete with a mechanism-level extraction: a risk taxonomy says what
can go wrong, a directory says who wrote about it, and neither records the
argument connecting a risk to a proposed remedy. Anchoring our extracted risks to
an external vocabulary of this kind is the natural next step
(\cref{sec:limitations}).

\paragraph{Causal and mechanism extraction.}
% [CITE: PICO extraction in evidence-based medicine; causal claim extraction.]
Evidence-based medicine has a mature tradition of extracting
population--intervention--comparison--outcome structure from trials. Our task is
the analogous reduction for a field whose evidence is largely theoretical, which
is why validation evidence appears as a reasoning stage rather than as an outcome
measurement.

\paragraph{Graph-based retrieval-augmented generation.}
% All IDs below VERIFIED against the source 2026-08-11. Entries in refs.bib.
GraphRAG systems build a graph over a corpus so that a language model retrieves
structured context rather than isolated passages
\citep{peng2024graphrag,sarmah2024hybridrag}; the design space has been mapped
modularly \citep{cao2024legographrag} and pushed towards learning graph
construction against downstream task reward \citep{tsang2025autographr1}.
\citet{pan2024unifying} survey the broader programme of joining language models to
knowledge graphs, and \citet{wu2025medgraphrag} build a medical GraphRAG whose
selling point is credible sourcing. In all of that work the graph is
\emph{instrumental}: it exists to improve an answer, and its edges are whatever
the indexing step produced. Our graph is the object of study, each edge is a
confidence-scored assertion that one identified paper made, and a query returns
that paper's own argument with provenance attached rather than a synthesised
answer. A GraphRAG layer over this dataset is an obvious application we do not
attempt here.

\paragraph{LLM extraction validation.}
We adopt the LLM-as-a-judge paradigm \citep{zheng2023judging} with an explicit
judge-of-judge calibration step, in which several independent stronger models
score the judge's proposals. \Cref{sec:limitations} states its scope limits: no
human adjudicated any verdict, so the validation chain is LLM-internal end to
end.

%%%%%%%%%%%%%%%%%%%%%%%%%%%%%%%%%%%%%%%%%%%%%%%%%%%%%%%%%%%%%%%%%%%%%%%%
\section{Methods}

\subsection{Corpus and Extraction}

We use the Alignment Research Dataset, a curated collection of AI safety
publications spanning arXiv preprints, LessWrong and Alignment Forum posts, and
organisational reports. After filtering to documents that yield a parseable
extraction, the working corpus is 11{,}779 documents.
% SRC: experiment_query_demo_report.json -> coverage_bound.papers_with_risk = 11779

Each document is passed in full to \texttt{o3} with a schema-constrained
extraction prompt (\cref{app:prompt}). The model traces the paper's own logical
chain and emits:

\begin{itemize}
\setlength{\itemsep}{0pt}
\item \textbf{Concept nodes}, each tagged with one \texttt{concept\_category}
from \{\emph{risk}, \emph{problem analysis}, \emph{theoretical insight},
\emph{design rationale}, \emph{implementation mechanism}, \emph{validation
evidence}\}.
\item \textbf{Intervention nodes}, each carrying
\texttt{intervention\_lifecycle} $\in 1..6$ (position in the research-to-deployment
cycle) and \texttt{intervention\_maturity} $\in 1..4$ (1 = conceptual proposal,
2 = early prototype or experiment, 3 = mature method, 4 = deployed in practice).
\item \textbf{Typed edges} carrying a relation type, a natural-language
description of the claimed relationship, and \texttt{edge\_confidence} $\in 1..5$
reflecting how explicitly the paper asserts the link.
\end{itemize}

The canonical chain the schema encourages is
\emph{risk} $\rightarrow$ \emph{problem analysis} $\rightarrow$ \emph{theoretical
insight} $\rightarrow$ \emph{design rationale} $\rightarrow$ \emph{implementation
mechanism} $\rightarrow$ \emph{validation evidence} $\rightarrow$
\emph{intervention}. Papers frequently deviate from it, omitting stages or
supplying several nodes for one stage, and we quantify that deviation in
\cref{sec:claim1} rather than forcing conformity.

Per-paper graphs are embedded (1536-dimensional OpenAI embeddings over node name,
aliases and description) and loaded into FalkorDB, giving 200{,}525 nodes:
19{,}096 risks, 36{,}959 interventions, and 144{,}470 intermediate concept nodes
(27{,}748 problem analysis, 26{,}361 theoretical insight, 27{,}543 design
rationale, 34{,}222 implementation mechanism, 28{,}596 validation evidence).
% SRC: experiment_race_prevalence_report.json -> by_node_type_full_graph.*.n
%      (n per type; 19096+27748+26361+27543+34222+28596+36959 = 200525)

\paragraph{Corpus composition.} ARD is heterogeneous and the chain dataset
inherits that heterogeneity: of the 2{,}772 chains, 1{,}055 come from arXiv, 326
from the Alignment Forum, 267 from LessWrong, 265 from the EA Forum, 238 from
YouTube transcripts, 120 from MIRI, and the remainder from a long tail of
institutional and personal sites. Seven chains (0.3\%, from 3 documents)
originate from Google Docs, at least one of which is a reading list rather than a
research document --- a reminder that ``document'' in ARD does not always mean
``paper''.
% SRC: recomputed 2026-08-11 over paths_hopwise_v4_edge_only_deduped.jsonl,
%      grouping by the source URL host of each chain's risk node.

\paragraph{Extraction is single-source.} Every structural edge comes from one
paper. The global graph concatenates per-paper graphs, and the pipeline never
asserts a relation between concepts introduced by two different papers. This
deliberate fidelity choice has consequences for every downstream graph statistic,
which we make explicit below and in \cref{sec:claim3}.

\paragraph{Extraction failures, and the limits of automated repair.}
Extraction fails outright on a minority of ARD records. The large majority of
those failures are records with no usable source text --- ARD entries carrying a
title and URL but an empty or \texttt{None} body --- which are unrecoverable in
principle. Filtering to failures repairable in principle (a URL, non-empty source
text, and at least one extracted node) leaves \textbf{441 candidates}. Running
the judge over all 441 in repair mode produced a non-empty graph for \textbf{23
of them (5.2\%)}, and those recovered graphs are thin: a mean of 2.4 nodes and
0.9 edges, against 16.0 nodes and 10.8 edges for audited successful extractions.
The practical conclusion, which we state because it is useful to others
attempting the same thing: \emph{an LLM judge audits and repairs an extraction
that already exists effectively, and reconstructs one that failed poorly.} We
therefore treat failed extractions as corpus loss rather than as a recoverable
pool.
% SRC: experiment_judge_full_report.json ->
%      item3_recovery.population_A_extraction_error_candidates (441 attempts, 23 with
%      a non-empty final_graph = 5.2%, mean 2.35 nodes / 0.87 edges)
% NOTE: the frequently-quoted "~60 recovered of ~400" figure is NOT used here. The 65
%      files in recovered_errors_graph/ are disjoint from these 441 attempts (overlap
%      = 0) and come from a different failure population whose denominator is not in
%      the bundle -- see item3_recovery.PROVENANCE_FINDING.

\subsection{Deduplication and Edge Normalisation}
% Reused near-verbatim from Gleb_overleaf_paper.txt L143-168.

The raw extraction produced duplicate nodes representing identical concepts with
different phrasings (``US semiconductor export controls'' and ``AI chip export
control regulations'') and inconsistent edge types (\textit{mitigated\_by} versus
\textit{mitigates}, or \textit{design\_rationale} versus
\textit{design\_rationale\_for}). We applied a two-threshold deduplication:
candidate pairs required both cosine similarity $\geq 0.88$ on embeddings and
Jaccard overlap $\geq 0.05$ on node names. We blocked by node type and category
to reduce comparisons, used FAISS \texttt{IndexIVFFlat} for approximate
similarity search, and applied Union-Find to group transitive duplicates. For
merged nodes we selected the canonical representative by longest description and
averaged embeddings. \Cref{app:dedup} reports four threshold configurations. We
prioritised precision over recall, because false merges corrupt downstream
pathway analysis by conflating distinct concepts whereas missed duplicates only
introduce redundancy. Edge normalisation unified bidirectional pairs
(\textit{mitigates} $\rightarrow$ \textit{mitigated\_by}) and standardised naming
variants, yielding 12 canonical edge types. Deduplication removed 2{,}385 nodes
(1.2\%).

\paragraph{Honest scope of the merge.} The merge is approximate in two respects
that we quantify rather than hide. First, the approximate nearest-neighbour
search under-retrieves: re-running the same threshold specification with
exhaustive search finds 54{,}282 candidate pairs against the 4{,}411 found by the
deployed configuration, a factor of 12.
% SRC: experiment_J6_merge_approx_v2_report.json -> variants.V2_T1_T2_T3.phase2_both_pass
%      = 54282 vs gleb_reference.candidate_pairs = 4411
Second, 1{,}140 nodes across 448 groups remain exact-name duplicates within their
category after the merge.
% SRC: experiment_J6_merge_approx_v2_report.json -> provenance_diagnostic
Both facts matter only where node identity is load-bearing. \Cref{sec:claim3}
shows exactly where that is, and the chain-level results in \cref{sec:claim1,sec:claim2}
do not depend on the merge.

\subsection{Structural Diagnostics of the Concatenated Graph}
% Reused from Gleb_overleaf_paper.txt L172-180, relocated from Results to Methods
% per paperA_positive_outline §3 ("never in Results").

The concatenated graph is highly fragmented: 15{,}123 disconnected components,
the largest containing 61 nodes (0.03\%), average degree 2.02, and an average
clustering coefficient of 0.013.
% SRC: re-derived from graph_node_attributes.pkl + graph_edge_data.pkl (EDGE type
%      only) 2026-08-11, union-find over 200,525 nodes / 202,149 EDGE edges.
%      NOTE: the frozen Overleaf reported 18,424 / 55 / 2.0 / 0.007 -- measured on
%      the MERGED 200,061-node graph with 197,542 edges, a different substrate.
%      We report our own numbers for the graph this paper describes.
This is not a finding about the literature. It is the arithmetic consequence of
concatenating per-paper chains without asserting cross-paper edges: a corpus of
$N$ short chains is exactly $N$ components until something links them. We record
it in Methods so that no later section reads as discovering it.

\subsection{Similarity Augmentation and Quality-Controlled Path Extraction}
\label{sec:paths}

Two graph layers are available for analysis, and the distinction governs every
claim in this paper.

\paragraph{EDGE layer (structural).} The extracted relations: single-paper,
LLM-asserted, confidence-scored.

\paragraph{SIM layer (similarity).} Cosine-similarity edges between nodes of the
same concept category, added to connect content across documents. At
$\tau \geq 0.80$ this adds 1{,}435{,}806 within-category edges, reduces components
from 15{,}123 to 4{,}124, and expands the largest component from 61 nodes to
152{,}753 (76.2\%). That expansion is the point, and also the hazard: a threshold
turning 0.03\% of the graph into 76.2\% of it in one step is not a neutral
operation (\cref{app:sensitivity}).
% SRC: re-derived 2026-08-11 from graph_edge_data.pkl, SIMILARITY edges with
%      cos_sim >= 0.80 restricted to same concept_category.
%      NOTE: the frozen Overleaf reported 169,083 edges -> 6,522 components ->
%      142,772 (71.4%); that k-NN construction is ~8.5x sparser than the dump this
%      paper uses. Same qualitative conclusion, different counts; we report ours.

For every quantitative result in \cref{sec:claim1,sec:claim2} we use the
\textbf{EDGE-only path set} under strict quality cuts applied at
path-construction time: structural edges only with no similarity hops, edge
confidence $\geq 3$, intervention-endpoint maturity $\geq 3$, and exactly one
risk node which must be the path root. Enumerating all such paths yields 8{,}954
raw chains.
% SRC: phase2_step4_F2v4_hopwise_falkordb.py (EDGE_CONFIDENCE_MIN=3,
%      INTERVENTION_MATURITY_MIN=3, single-risk); verified empirically in
%      experiment_paper_claim_audit.json: 100% of paths start at a risk node, 0 paths
%      contain a risk after position 0, all endpoints have maturity 3 (8050) or 4 (904).

\paragraph{The reporting unit.} Raw path enumeration over-counts prolific papers:
one paper contributes 700 of the 8{,}954 raw chains (7.8\%), two others
contribute 200 and 184, and 169 papers contribute ten or more. Reporting
per-chain statistics on that set would weight the corpus by how branched each
paper's extracted graph happens to be --- the same unit-of-analysis error we
document in \cref{sec:claim3}. We therefore collapse near-duplicate chains within
each source paper: two chains are near-duplicates when the smaller one's node set
is $\geq 70\%$ contained in the larger's, and the shorter is dropped. This yields
the \textbf{2{,}772-chain de-duplicated set} (at most 7 chains from any one
paper), the reporting unit for every number below. Both sets cover the same
\textbf{1{,}868 source papers}, and we report the raw set alongside as a
sensitivity check.
% SRC: phase1_dedup_paths.py (CONTAINMENT_THRESHOLD = 0.70 -- note the module
%      docstring says 80%, the code uses 70%; the code is what ran).
%      Counts + skew: experiment_paper_claim_audit.json -> path_set_comparison.

The path set is free of similarity-threshold sensitivity and of merge artefacts
by construction. It is the substrate on which we are willing to make numerical
claims.

\subsection{Validation}
\label{sec:validation}

We assess extraction quality with an LLM-as-a-judge protocol. A judge model from
a different provider than the extractor (\texttt{claude-sonnet-4-5}, batch API)
receives the source text together with the extracted graph and returns a
structured audit: schema conformance, referential integrity, orphan and duplicate
detection, rationale consistency, coverage of the source's claims, and a set of
proposed repairs. We ran it on a random sample of \textbf{100 successfully
extracted papers} spanning all nine ARD source types (LessWrong 28, Alignment
Forum 21, blogs 15, EA Forum 11, arXiv 11, Arbital 6, special documents 4,
YouTube transcripts 2, aisafety.info 2). A judge-of-judge stage then had three
independent meta-graders (\texttt{claude-opus-4-5}, \texttt{gemini-3-pro},
\texttt{gpt-5.1}) score each extraction before and after the judge's proposed
repairs.
% SRC: experiment_judge_item2_report.json (n=100, by_source_type)
%      + extraction_validator/STATUS.md on branch judge_handoff_workshop_items_2_3

\paragraph{Structural error classes.} Across the 100 papers the judge found
\textbf{108 referential-integrity findings}, of which 68 are blocker-severity and
the dominant pattern is an edge whose source or target names a node absent from
the node list; \textbf{42 orphan nodes} (present in 40\% of papers); and
\textbf{56 within-paper duplicate node pairs} (55\% of papers). The mean audited
graph has 16.0 nodes and 10.8 edges.
% SRC: experiment_judge_item2_report.json -> findings_total,
%      findings_papers_with_at_least_one, per_paper_means

\paragraph{Coverage.} Asked to identify content in the source that the extraction
missed, the judge proposed \textbf{new nodes for only 7 of 100 papers} (9 nodes
in total). Its coverage list flags a mean of 7.8 relationships per paper that it
judges the source to support but the extraction not to contain --- an upper-bound
opinion signal rather than a measured omission rate, since the judge has no
schema slot in which to propose those edges and no human adjudicated them.
% SRC: experiment_judge_item2_report.json -> proposed_fix_totals.add_nodes = 9;
%      findings_papers_with_at_least_one.add_nodes = 7; per_paper_means.missing_edges_flagged

\paragraph{Judge-of-judge.} Three meta-graders scored each extraction on a 0--100
holistic rubric before and after applying the judge's proposed repairs
(\cref{tab:graders}). All three agree on direction: every grader records an
improvement on every paper it scored, and on the 13 papers all three scored,
agreement that the judge improved the extraction is unanimous. They disagree
sharply on magnitude, with mean gains from $+3.5$ to $+33.1$ points.
% SRC: experiment_judge_full_report.json -> item2_meta_graders,
%      item2_grader_agreement (reproduces extraction_validator/results.md exactly)

\begin{table}[t]
\centering
\small
\begin{tabular}{lrrrr}
\toprule
\textbf{Meta-grader} & \textbf{$n$} & \textbf{pre} & \textbf{post} & \textbf{$\Delta$} \\
\midrule
Claude Opus 4.5 & 95 & 69.7$\pm$18.0 & 73.1$\pm$18.2 & $+3.5$ \\
GPT-5.1         & 95 & 67.0$\pm$16.2 & 80.2$\pm$10.2 & $+13.2$ \\
Gemini 3 Pro    & 13 & 62.7$\pm$22.6 & 95.8$\pm$1.8  & $+33.1$ \\
\bottomrule
\end{tabular}
\caption{Meta-grader rubric scores before and after the judge's repairs. The
three graders scored \emph{unequal, partly disjoint} subsets of the 100 papers
--- only 13 papers carry all three scores --- because the rubric ran across
several prompt iterations and only some iterations emit a paired pre/post score.
The rows are therefore not directly comparable, and Gemini's $n{=}13$ post-repair
mean of 95.8 with a standard deviation of 1.8 shows saturation at the top of the
scale.}
\label{tab:graders}
\end{table}

Inter-grader agreement on the 13 commonly scored papers is moderate before the
repairs (Fleiss' $\kappa = 0.54$ on four ordinal score bands) and collapses
afterwards ($\kappa = 0.09$): Opus and GPT-5.1 remain rank-correlated on the
post-repair scores (Spearman $\rho = 0.50$) while Gemini's saturated scores
correlate negatively with both ($\rho = -0.63$, $-0.64$). We therefore report the
\emph{direction} of the judge's effect as a robust, unanimous finding and treat
its \emph{magnitude} as unresolved.
% SRC: experiment_judge_full_report.json -> item2_grader_agreement

\paragraph{What the post-repair score is.} The meta-graders saw the source text,
the extraction, and the judge's report including its \emph{proposed} fixes, and
scored the extraction before and after those fixes. The post-repair score
therefore assesses the proposed repairs as described rather than measuring a
re-built graph. This is a real limitation, since no repaired graph was
re-extracted and re-scored. It also means the known defects in the fix-application
code, which under-applied multi-node additions and dropped proposed merges, do
\emph{not} contaminate these scores.
% SRC: prompts/judge_and_extraction_evaluation_prompt.md ("give a score between
%      0-100 before and after the judge's proposed fixes") + combine_judge_and_
%      extraction_with_original_text.py, which bundles Original Text + Extraction
%      Output + Judge Output for the meta-grader. Verified 2026-08-11; this
%      corrects an earlier working assumption that the graders scored the applied
%      final_graph.

\textbf{[GAP: a human-anchored spot-check. No human adjudicated any judge or
meta-grader verdict, so the whole validation chain is LLM-internal. Owner: Sai,
ticket \#150, together with the manual 50-instance error taxonomy.]}

\paragraph{Two measurement caveats we do not hide.} First, the judge grades
against a schema expecting rationale fields inline on nodes and edges, whereas
the pipeline stores rationale as separate nodes linked by a dedicated relation.
That mismatch accounts for \textbf{184 of the 218 blocker-severity schema flags
(84\%)}; excluding it leaves 34 substantive blocker-severity schema findings
across 100 papers. Any headline error rate computed from raw judge output without
this split measures a schema version difference, not extraction quality.
% SRC: experiment_judge_item2_report.json -> schema_check_severity_rationale_field_mismatch
%      vs schema_check_severity_substantive
Second, the judge scores node and edge faithfulness, referential structure and
edge type. It does \emph{not} score the \texttt{intervention\_maturity}
assignment. Wherever maturity appears below we label it LLM-assessed and use it
descriptively, never as a validated measurement.

%%%%%%%%%%%%%%%%%%%%%%%%%%%%%%%%%%%%%%%%%%%%%%%%%%%%%%%%%%%%%%%%%%%%%%%%
\section{C1: The Chains Carry Argument Structure}
\label{sec:claim1}

Every number in this section is a structural property of the EDGE-only
quality-cut path set (\cref{sec:paths}) or a direct count over extracted
attributes. No clustering, no centrality, and no similarity threshold enters any
of them.

\subsection{Chains carry the full reasoning structure}

Across the 2{,}772 de-duplicated chains, \textbf{87.4\% (2{,}423) contain all
five intermediate reasoning stages}. Stage-wise presence is: problem analysis
99.0\%, theoretical insight 93.9\%, design rationale 96.1\%, implementation
mechanism 96.0\%, validation evidence 98.0\%. On the raw 8{,}954-chain set the
same figure is 89.0\%, so the result does not depend on the de-duplication
choice.
% SRC: experiment_paper_claim_audit.json -> path_set_comparison.deduped
%      (all5 2423/2772 = 87.4%; subtype_presence_pct) and .raw (89.0%)

A large majority of extracted chains therefore do more than connect a risk to an
intervention: they carry the paper's argument for why the intervention should
work, stage by stage.

\paragraph{What this is not.} 87.4\% is the share \emph{of complete chains}
carrying all five stages, not the share of the corpus. Only 1{,}868 of the
11{,}779 documents (15.9\%) yield any complete high-confidence chain; the
remaining documents contribute nodes and partial structure. Both numbers matter:
the corpus-level yield is modest, and what is yielded is structurally rich.

\subsection{The chains are not a filled-in template}

The obvious objection is that a schema suggesting a seven-node chain will produce
seven-node chains regardless of the source. \Cref{fig:dataset}A refutes the
strong form of that objection: chains span 4--16 nodes, the mode is 7 (1{,}775
chains, 64.0\%), but 25.8\% run \emph{longer} than the canonical seven --- carrying
multiple nodes for one or more stages, that is, richer argumentation --- and
10.2\% run shorter, carrying genuine gaps.
% SRC: experiment_paper_claim_audit.json -> path_set_comparison.deduped.len_hist
%      (=7: 1775 = 64.0%; >7: 715 = 25.8%; <7: 282 = 10.2%)

\paragraph{What the length spread does and does not settle.} A spread is
consistent with the extraction following each paper's argument, and it is also
consistent with a schema that bends a little under pressure. Settling the
question requires separating two failure modes that are easy to conflate.
\emph{Scope} failure means the pipeline emits a chain from a document that is not
about AI safety. \emph{Fidelity} failure means the pipeline emits a chain whose
structure the source document does not contain. Only the second bears on C1.

An out-of-domain control does not test fidelity. Any literature that names a
problem, analyses it, and proposes a remedy --- materials science, epidemiology,
software engineering --- contains genuine risk-to-intervention arguments, so a
high chain-emission rate on non-safety documents would show the schema is
domain-general, which is a scope property and arguably a feature. Two controls do
test fidelity, and both hold the domain fixed:

\begin{itemize}
\setlength{\itemsep}{1pt}
\item \textbf{Schema ablation.} Re-extract a sample with a prompt that asks for
the chain from problem to proposed solution \emph{without naming the five
stages}, then measure how often the emergent chain maps onto them. Structure that
survives un-prompted is read from the document rather than imposed on it.
\item \textbf{Degraded-source control.} Re-extract from versions of the same
safety documents whose argument has been destroyed but whose vocabulary is
retained --- sentence-shuffled text, abstract only, reference list only. A
complete five-stage chain arising at a similar rate would show the model
confabulating structure from topical vocabulary.
\end{itemize}

We ran neither, and we flag them as the most valuable additions to this work.
\Cref{app:failure} is the one fidelity failure we exhibit: a chain that passes
every structural filter, extracted from an expository article proving that there
are infinitely many primes, in which the model invents a safety framing the
source never asserts. We report it as an existence proof, not as a rate.

\paragraph{What C1 does not claim.} C1 concerns fidelity to what a paper argued.
It says nothing about whether the argument is \emph{correct}, whether the
intervention works, or whether the analysis is good safety research. A faithfully
extracted chain from a weak paper is a successful extraction.

\subsection{What the judge study supports}

\Cref{sec:validation} reports the judge results in full. Their bearing on C1 is
narrow and worth stating plainly. The judge proposed new nodes for 7 of 100
papers, and the Opus meta-grader recorded 216 missed concepts across the 95
papers it scored (at least one in every paper) against 20 instances of fabricated
content in 16\% of papers and 15 category errors in 33\%.
% SRC: experiment_judge_full_report.json -> item3_error_taxonomy_auto
Under this grader the dominant failure mode is \emph{omission, not
hallucination}. That supports reading an extracted chain as a partial record of
what a paper argued, and does not support reading the absence of a stage as
evidence that the paper lacked one. These counts are un-adjudicated LLM output
and we report them as a profile rather than as a measured error rate; the manual
taxonomy remains outstanding.

%%%%%%%%%%%%%%%%%%%%%%%%%%%%%%%%%%%%%%%%%%%%%%%%%%%%%%%%%%%%%%%%%%%%%%%%
\section{C2: Chains as the Unit of Analysis}
\label{sec:claim2}

\subsection{Mechanism-level retrieval}

The dataset supports a query that document-level resources cannot serve:
\emph{by what mechanism does intervention $I$ reduce risk $R$?} Because the
mechanism is stored as an explicit path, answering it is a single graph lookup
returning the chain and its source citation, with no clustering, no ranking, and
no embedding search over abstracts. \Cref{tab:query} shows two chains retrieved
this way, one governance and one technical, both from arXiv sources so that the
extraction input is structured and the citation is unambiguous.
% SRC: experiment_query_demo_report.json -> query_demo.*

\begin{table*}[t]
\centering
\small
\begin{tabular}{p{0.09\textwidth}p{0.86\textwidth}}
\toprule
\textbf{Stage} & \textbf{Query 1 --- governance.} Source: \citet{shavit2023chinchilla} \\
\midrule
risk & Covert development of dangerous ML models using large-scale training \\
pa & Tampering with logging mechanisms via supply chain attacks \\
ti & Open-source Hardware Roots-of-Trust enhance transparency and tamper evidence \\
dr & Embed open-source HRoT like Caliptra in ML chips \\
im & Signed firmware enforced by HRoT with anti-rollback protections \\
va & Industry adoption of Caliptra HRoT across major vendors \\
\textbf{intv} & Adopt open-source HRoT in data-center ML chips to enforce secure logging \textbf{(maturity 3)} \\
\midrule
\multicolumn{2}{c}{} \\
\midrule
\textbf{Stage} & \textbf{Query 2 --- technical alignment.} Source: \citet{ngo2022alignment} \\
\midrule
risk & Catastrophic disempowerment of humanity by misaligned AGI \\
pa & Misaligned goal formation during AGI training \\
ti & Reward misspecification in real-world RL tasks \\
dr & Automated scalable supervision to reduce reward misspecification \\
im & Reinforcement learning from human feedback scaling \\
va & Instruction-following improvement in InstructGPT via RLHF \\
\textbf{intv} & Scale RLHF fine-tuning using AI-generated feedback signals \textbf{(maturity 3)} \\
\bottomrule
\end{tabular}
\caption{Two mechanism-level retrievals, one governance and one technical, both
from arXiv sources. Each block is a single extracted chain returned verbatim by a
risk--intervention query; stage abbreviations follow the logical chain order
(pa = problem analysis, ti = theoretical insight, dr = design rationale,
im = implementation mechanism, va = validation evidence). Neither answer is
synthesised: each is one stored path with its source document attached.}
\label{tab:query}
\end{table*}
% SRC: Query 1 = node ids [1067,1086,1087,1088,1089,1090,1091], arxiv.org/abs/2303.11341
%      Query 2 = node ids [415,416,417,426,431,437,442],      arxiv.org/abs/2209.00626
%      Both verified against paths_hopwise_v4_edge_only_deduped.jsonl 2026-08-11.
%      Titles/authors verified against arXiv the same day.
%      These replace the earlier three examples: the original Query 1 was sourced from a
%      Google-Doc reading list, and the other two from non-arXiv sources. See
%      paperA_v2_AUDIT.md P3.

Three properties of these results matter. They are \emph{single-hit}: the answer
is one stored path, not a synthesis over retrieved passages. They are
\emph{attributable}: every node carries its source document. And they are
\emph{comparable}: because all chains share the same stage vocabulary, chains
from different papers addressing the same risk sit side by side, which is the
precondition for the cross-paper mechanism analysis we identify as the natural
next step (\cref{sec:limitations}).

\subsection{Reading the corpus at chain level}

The same unit supports descriptives that no metadata directory can produce,
because they are properties of the argument inside each paper.

\paragraph{Intervention maturity.} Of the 36{,}959 extracted interventions, the
LLM-assessed maturity distribution is \textbf{27.6\% conceptual} (10{,}189),
\textbf{57.3\% early-stage} (21{,}180), \textbf{12.7\% mature} (4{,}681), and
\textbf{2.5\% deployed} (909) (\cref{fig:dataset}B). Roughly six in seven proposed
interventions in this corpus sit at conceptual or early-stage maturity, and fewer
than one in thirty is described as deployed.
% SRC: experiment_dataset_strength_report.json -> interventions.maturity

\paragraph{How to read maturity.} Maturity is an LLM assignment that no human
adjudicated and that our judge study does not score (\cref{sec:validation}). We
read it as coarse composition rather than as measurement. The skew is a
four-bucket call with a very large margin and would survive a substantial
per-item error rate, but the individual bucket boundaries should not be read as
precise.
% NOTE 2026-08-11: a scope-composition paragraph (share of chains whose risk is a
%      capability gap rather than a human harm) was drafted here and WITHDRAWN. It
%      depended on phase2_routing_assignments.jsonl, which is not committed and is
%      not re-derivable by a reader at reasonable cost, and it belongs to a separate
%      line of work. Do not reinstate it unless that file ships with the release and
%      a committed script reproduces it. See paperA_v3_RESTRUCTURE.md §2.2.

\subsection{End-to-end coverage}

All 11{,}779 corpus documents yield at least one risk node and at least one
intervention node --- the two sets are identical --- which is a property of the
extraction schema, not a measurement of the literature. The informative number is
how many documents yield a \emph{complete, high-confidence, mature-intervention
chain end to end}: \textbf{1{,}868 documents} (15.9\%), contributing
\textbf{2{,}643 distinct risk$\rightarrow$intervention node pairs} after
de-duplication (3{,}222 before).
% SRC: experiment_paper_claim_audit.json -> n_source_papers = 1868 recomputed from
%      the raw path file (identical on the de-duplicated set);
%      distinct_ri_node_pairs deduped 2643 / raw 3222.
%      NOTE: experiment_query_demo_report.json carries the field
%      "papers_with_complete_strict_chain": 1868 but it is a HARDCODED literal in
%      experiment_query_demo.py:125, not computed there. Cite the audit receipt.
The ratio between the two figures is the honest statement of yield: roughly one
document in six argues a complete mechanism from a named risk to an intervention
that is at least a mature method, under our quality cuts.

%%%%%%%%%%%%%%%%%%%%%%%%%%%%%%%%%%%%%%%%%%%%%%%%%%%%%%%%%%%%%%%%%%%%%%%%
\section{C3: The Aggregate Graph Does Not Support Graph-Analytic Inference}
\label{sec:claim3}

A graph of this shape invites standard network analysis: cluster the nodes, rank
by centrality, compare risk and intervention densities per cluster, look for
co-occurring framings. We ran that analysis. This section reports what we
measured when we checked those quantities against construction artefacts. Each
failure mode is stated as a measurement on this corpus together with the control
that removes it.

\paragraph{Provenance of this section.} The analyses below are \emph{post-hoc}.
An earlier pass over this corpus produced a set of graph-analytic findings; we
then re-derived them under controls, and two did not survive. We report the
re-derivations rather than the originals, and we flag the post-hoc status because
a reader should weight a hypothesis formed after seeing the data less than one
registered before it.

\subsection{Near-duplicate merging manufactures a centrality super-hub}

The corpus contains thousands of near-identical existential-risk concept nodes
(``Existential catastrophe from misaligned advanced AI'' and close paraphrases)
contributed by different papers. Merging them is the natural cleanup step, and it
is exactly what breaks centrality, because the canonical node inherits every
member's structural edges. In our measurement the largest merged risk node
absorbed \textbf{4{,}066 members and 7{,}777 structural edges}.
% SRC: experiment_merge_vs_simexcl_ec_report.json -> biggest_merged_node_members / _edge_degree

\Cref{tab:ec} shows eigenvector centrality under the four combinations of
(merged / un-merged) $\times$ (similarity edges included / risk--risk similarity
excluded).

\begin{table}[t]
\centering
\small
\begin{tabular}{lrrr}
\toprule
\textbf{Condition} & \textbf{EC$_1$} & \textbf{EC$_1$/EC$_2$} & \textbf{xrisk} \\
                   &                 &                        & \textbf{in top-10} \\
\midrule
un-merged, full SIM        & 0.0328 & 1.0$\times$ & 10/10 \\
un-merged, risk--risk SIM excl. & 0.0011 & 1.0$\times$ & 8/10 \\
merged, full SIM           & 0.5069 & 33.6$\times$ & 1/10 \\
merged, risk--risk SIM excl.    & 0.3337 & \textbf{90.1$\times$} & 2/10 \\
\bottomrule
\end{tabular}
\caption{Eigenvector centrality of the top node under four graph-construction
choices. Merging near-duplicate risk nodes produces a single node with 90$\times$
the centrality of the next node, and excluding similarity edges does not repair
it. On the un-merged graph the top nodes fall within 1\% of each other: the
corpus contains no dominant hub, only one manufactured by the merge.}
\label{tab:ec}
\end{table}
% SRC: experiment_merge_vs_simexcl_ec_report.json -> conditions[0..3].top10[0].ec
%      and ratio to top10[1].ec : 0.032753/0.032686=1.00 ; 0.0010682/0.0010393=1.03 ;
%      0.50685/0.015084=33.6 ; 0.33369/0.0037034=90.1

Two consequences follow, and their order matters. \textbf{(a)~Un-merging is
primary.} On the merged graph the super-hub retains a 90$\times$ lead even after
deleting all risk--risk similarity edges, because the merge baked the
concentration into the structural topology. \textbf{(b)~Excluding risk--risk
similarity is still necessary}, but as the \emph{second} step: on the un-merged
graph it flattens the top centrality value by a factor of 31 (0.0328
$\rightarrow$ 0.0011) by breaking the near-duplicate similarity clique. Applied
alone, either control leaves the artefact intact.

The substantive point for readers of this literature: \emph{``the most central
risks are existential'' is a statement about graph construction, not about the
field.} On the un-merged graph the ten highest-centrality nodes differ from each
other by less than 1\%, so no ranking exists to report.

\paragraph{Control.} Do not merge near-duplicates before computing centrality.
Exclude within-category similarity edges from all centrality measures. Count
canonical concepts or source papers, never raw nodes, for any co-occurrence
statistic.

\subsection{Approximate near-duplicate search under-removes, unpredictably}

Deduplication quality is usually reported as a removal rate. That number is
uninformative on its own: re-running our own threshold specification with
exhaustive rather than approximate search returns \textbf{54{,}282 candidate
pairs against 4{,}411}, a factor of 12, and the gap concentrates almost entirely
in the risk block, precisely where the centrality artefact lives.
% SRC: experiment_J6_merge_approx_v2_report.json -> variants.V2_T1_T2_T3 vs gleb_reference

\paragraph{Control.} Report recall of the candidate-generation stage, not only
the removal rate. Verify recall on the node block that downstream metrics are
most sensitive to.

\subsection{Loose similarity thresholds group topics, not mechanisms}

Direct $k$-means on the 1536-dimensional embeddings fails outright (silhouette
$\approx 0.02$). UMAP reduction to 150 dimensions followed by $k$-means at
$k{=}40$ reaches silhouette 0.298 --- modest against conventional interpretation
thresholds \citep{kaufman1990finding} (\cref{app:clustering}) --- and the
resulting clusters are recognisable field topics (\cref{app:clusters}). But
cluster co-membership at $\tau{=}0.80$ tracks vocabulary rather than mechanism:
the same threshold producing navigable topics also grows the largest connected
component from 61 nodes to 152{,}753. We attempted to derive a clean,
non-overlapping risk/intervention taxonomy this way, tightening quality cuts step
by step --- maturity, edge confidence, EDGE-only paths, clustering each node type
on its own axis --- and node-level semantic clustering remained insufficient to
carry mechanism-level conclusions.

\paragraph{Control.} Treat node-level clusters as a browsing index, not as an
inferential unit. Cluster risk and intervention nodes on separate axes. Do not
compute risk-to-intervention ratios per cluster, because numerator and
denominator are populated by different extraction stages and the ratio inherits
the merge artefact above.

\subsection{Co-occurrence measured on a centrality-selected head is circular}

Selecting the top-100 risks by eigenvector centrality and measuring which framings
appear there is a natural exploratory move, and in an earlier pass on this corpus
it produced a striking-looking result: 88\% of the single-path risks in that head
appeared to be framed through competitive race dynamics, with an apparent
44-fold gradient from the centrality head to the full risk population. Neither
figure survives re-derivation. Re-running the same measurement under the same
definitions --- importance by eigenvector centrality, path diversity by distinct
first-hop structural problem-analysis neighbours, race framing by the same
keyword rule --- reproduces the \emph{structural} part (38 of the top-100 risks
are single-path, against 41 reported) but not the framing part: 2.6\% of those
single-path risks are race-framed, and the head-to-population gradient is
2.0-fold, not 44-fold. The keyword classifier itself is sound ---
\cref{app:race} reports 92\% precision on a stratified manual sample --- so the
failure lies in the selection step, not in the measurement instrument.
% SRC: experiment_race_top100_rederive_report.json -> condition_A (merged + full
%      within-category SIM, the closest reconstruction of the original setup).
%      The similarity layer in the graph dump we use is ~8.5x denser than the one
%      the original analysis ran on, so the centrality ranking -- and hence the
%      selected head -- is not identical. We report the failure to reproduce rather
%      than the original figure.
Corpus-wide prevalence of the framing is \textbf{2.21\% of problem-analysis
nodes} (614/27{,}748) and 1.65\% of risk nodes; only 0.50\% of EDGE-only chains
contain a race-framed problem-analysis node, and the re-derived head rates
(1.3--6.3\% across centrality tiers) sit within a small multiple of that base.
% SRC: experiment_race_prevalence_report.json + experiment_race_top100_rederive_report.json

\paragraph{The lesson is stronger, not weaker, for the reproduction failing.} A
selection-conditioned statistic on this kind of graph is not merely inflated, it
is \emph{unstable}: it depends on the similarity-layer density and the merge
decisions that set the ranking, and those are pipeline parameters rather than
properties of the literature. Nobody should report a number computed this way
without stating the selection rule, the graph construction, and a full-population
baseline alongside it.

The underlying association is real but modest. \Cref{tab:race} shows the
race$\times$existential-risk co-occurrence odds ratio under three units of
analysis. It survives collapsing near-duplicates, which rules out pure
pseudo-replication, but at 2.7--5.2$\times$ on a $\sim$2\% base rate it supports
no claim about the field's discourse structure.

\begin{table}[t]
\centering
\small
\begin{tabular}{lrrr}
\toprule
\textbf{Unit} & \textbf{OR} & \textbf{95\% CI} & \textbf{$P$(race)} \\
\midrule
node               & 3.82 & 3.24--4.50 & 7.6\% vs 2.1\% \\
merged node        & 2.73 & 2.06--3.64 & 5.5\% vs 2.1\% \\
source paper       & 5.17 & 4.32--6.20 & 10.7\% vs 2.3\% \\
\bottomrule
\end{tabular}
\caption{Race-framing $\times$ existential-risk co-occurrence across units of
analysis (all $p < 10^{-12}$). $P$(race) columns give the rate within
existential-risk versus non-existential-risk nodes. The association is
directionally robust and quantitatively minor.}
\label{tab:race}
\end{table}
% SRC: experiment_race_dedup_robustness_report.json -> A_node_level / B_merged_node_level / C_paper_level

A related trap: deleting race-framed problem-analysis nodes appeared to
disconnect 51 of the top-100 risks from all interventions. Recomputed on the
corrected merge, the same deletion isolates \textbf{2 of 100}.
% SRC: experiment_J6_merge_approx_v2_report.json -> j6_on_V2.isolated_after_race_removal = 2
%      vs gleb_reference.j6_isolated = 51

\paragraph{Control.} Report prevalence on the full population before reporting
any pattern in a selected head. State the selection rule as part of the result.
Choose the unit of analysis (node, merged node, paper) explicitly and show that
the effect survives all three.

\subsection{Single-source extraction makes several ``findings'' true by design}

Three quantities that look like discoveries are consequences of the pipeline.
\emph{Path diversity}: because each paper contributes one chain, a risk discussed
in one paper necessarily has one route to an intervention, which is extraction
structure rather than fragility in the literature. \emph{Path completion}: under
risk-rooted breadth-first search over a similarity-augmented graph most branches
terminate at intermediate stages, so a low completion percentage measures the
traversal, not the corpus. \emph{Siloing}: cross-paper links exist only in the
similarity layer, so the structural layer is siloed by construction while the
similarity layer's connectivity is a function of $\tau$, and no layer measures
real siloing.

\paragraph{Control.} Before reporting any structural statistic, ask what value it
would take on a corpus of independent single-paper chains. If the answer is the
value you observed, it is a pipeline property.

%%%%%%%%%%%%%%%%%%%%%%%%%%%%%%%%%%%%%%%%%%%%%%%%%%%%%%%%%%%%%%%%%%%%%%%%
\section{Reproducibility and Release}
\label{sec:repro}

Every quantitative claim in this paper carries a source comment in the manuscript
naming the receipt file it came from, and one script re-derives all of them from
the raw path files and the node attributes, then compares each against the value
printed here. On the current draft it reports 42 of 42 numeric claims passing. A claim-to-script-to-receipt map ships with the code.

The extraction pipeline, the path builders, the audit scripts and the receipt
JSONs are released together with the chain dataset under a permissive licence.
Two large intermediate checkpoints (node attributes and edge data, together
$\sim$3.2 GB) are rebuildable from the released graph dump with a released
script; because the receipts are committed, every number here is checkable
without them. \Cref{fig:dataset} plots only from committed receipts for the same
reason.

\paragraph{Compute.} Extraction used the OpenAI batch API with \texttt{o3} over
11{,}779 documents. The judge and meta-grader runs used the Anthropic, OpenAI and
Google batch APIs at the model versions named in \cref{sec:validation}; model
identifiers are pinned rather than floating, so a re-run reproduces the same model
generation. All graph analysis runs on one workstation CPU without a GPU, with
path enumeration against the graph database as the heaviest step.
% [GAP: exact wall-clock and token totals for the extraction run, if recoverable from
%  the batch-API logs. Not blocking.]

\subsection{Use of AI assistance}
\label{sec:aidisclosure}

Beyond the extraction and judging models, which are the object of study and are
described in \cref{sec:validation}, AI tools were used in preparing this work. An
LLM coding assistant (Anthropic Claude) wrote and ran the analysis and audit
scripts released with the paper, and drafted and restructured portions of this
manuscript. All claims, all numerical results and all interpretation are the
authors' own: every quantitative statement is re-derived from the raw data by the
released audit script and was checked against it, and the authors take full
responsibility for the content. No LLM is an author, and none is credited as one.
% [GAP: scope of the drafting claim. "portions of this manuscript" is deliberately
%  imprecise until the co-authors confirm what each of them wrote by hand -- parts of
%  Methods and several appendices are reused near-verbatim from a co-author's frozen
%  Overleaf draft, and this statement must not overstate the assistant's share or
%  understate theirs. Settle alongside the author list and CRediT statement (G15);
%  the same venue policies that require this disclosure also bar an LLM from
%  authorship, so the two decisions belong in one conversation.
%  Re-check the wording against the chosen venue: ICLR 2026 desk-rejects undisclosed
%  LLM use and requires all contributions to be acknowledged; ICML 2026 permits
%  assistance but places all responsibility on the authors and forbids crediting an
%  LLM. Verified 2026-08-11.]

%%%%%%%%%%%%%%%%%%%%%%%%%%%%%%%%%%%%%%%%%%%%%%%%%%%%%%%%%%%%%%%%%%%%%%%%
\section{Limitations and Outlook}
\label{sec:limitations}

\paragraph{No comparative baseline.} We claim that this extraction recovers
argument structure, not that it beats an alternative extractor. We compare
against no simpler pipeline --- flat triple extraction, abstract-only extraction,
a non-reasoning model --- and we ran no ablation over schema components.
A reader should therefore read C1 as evidence that the output has the claimed
structure, not as evidence that the design choices are necessary for it. The
schema-ablation and degraded-source controls described in \cref{sec:claim1} are
the highest-value missing experiments, and the first doubles as the missing
ablation.

\paragraph{Corpus.} ARD is a curated snapshot and its coverage thins after 2023,
so frontier-era discourse is under-represented. Claims here describe the extracted
corpus, not the field as of publication.

\paragraph{Single extractor.} All extractions come from one model (\texttt{o3})
under one prompt. Our validation design does not characterise model-specific
extraction bias.
% [GAP: Mike ran an n=20 multi-model consistency check (o3 / GPT-5 / Claude-4).
%  If those numbers are recoverable, one sentence here materially strengthens this
%  limitation. Owner: Sai/Mike via ticket #150 recovery bundle.]

\paragraph{Validation is LLM-based.} An LLM judges faithfulness with
judge-of-judge calibration; domain experts do not, at any scale. No human
adjudicated a single verdict in the chain.

\paragraph{LLM-assigned attributes.} Concept subtype, edge confidence,
intervention lifecycle and maturity are model assignments. The judge scores only
node and edge faithfulness and edge type.

\paragraph{Dual-use and misreading risk.} The most likely harm from this dataset
is misuse of its aggregate statistics: a reader who takes cluster sizes or
centrality rankings from the released graph and reports them as facts about which
risks the field neglects will produce a confident and wrong claim.
\Cref{sec:claim3} exists partly to make that harder, and the release
(\cref{sec:repro}) carries the same warning.

\paragraph{Outlook.} The natural next step, and the reason for the release, is
cross-paper mechanism analysis: grouping interventions into mechanism families and
anchoring risks to an external risk taxonomy, so that one can ask which mechanism
families address which risks, at what stage of validation evidence, across the
whole corpus --- and, where a mechanism family is densely populated for every risk
but one, whether that gap marks a research opportunity. We deliberately do not
attempt it here, because \cref{sec:claim3} shows that node-level semantic
clustering on this data is too weak to carry mechanism-family grouping, and doing
it properly requires an external anchor rather than a better threshold.

%%%%%%%%%%%%%%%%%%%%%%%%%%%%%%%%%%%%%%%%%%%%%%%%%%%%%%%%%%%%%%%%%%%%%%%%
\section{Conclusion}

We extracted, from 11{,}779 AI safety documents, the explicit logical chains by
which proposed interventions are argued to reduce named risks. Across 2{,}772
quality-controlled chains, 87.4\% carry all five reasoning stages while spanning
4--16 nodes, structure captured from the papers rather than imposed on them. The
chains answer mechanism-level queries that document-level resources cannot, and
they expose what the corpus holds: 2.5\% of extracted interventions are described
as deployed, and one document in six yields a complete chain at all. We also
report, as a first-class result, four measured ways in which
routine graph analysis of such a corpus produces artefacts rather than findings
--- most sharply a merge-manufactured centrality hub at 90$\times$ the next node
--- together with the controls that remove them, and two previously reported
findings on this corpus that fail to reproduce. The chain is the unit that
survives scrutiny on this data; the aggregate graph is not.

%%%%%%%%%%%%%%%%%%%%%%%%%%%%%%%%%%%%%%%%%%%%%%%%%%%%%%%%%%%%%%%%%%%%%%%%
\section*{Acknowledgments}

This work was carried out by a distributed volunteer team convened through
\textbf{EleutherAI}, whose Discord server hosted all project coordination ---
stand-ups, working threads and review discussion --- for the full duration of the
project. We thank EleutherAI for providing and maintaining that platform.

% 🔴 GATE -- MUST BE CLOSED BEFORE SUBMISSION. Do not remove this comment until resolved.
%
% COMPUTE ACKNOWLEDGEMENT. Model inference for the extraction and judge runs was paid
% for with API credits provided by a private acquaintance. Before this paper is
% submitted, Martin must confirm with them whether they wish to be acknowledged by
% name, acknowledged anonymously, or not mentioned at all.
%
% Why this is a gate and not a nicety: the credits may have been granted under terms
% that did not contemplate being passed to a third party. Naming the donor without
% asking could expose them to a problem with whoever issued the credits. Naming an
% institution that did not authorise the transfer would be worse.
%
% Three acceptable outcomes, pick one and delete the other two:
%   (a) named:     "Compute for this work was donated by <NAME>."
%   (b) anonymous: "Compute for this work was donated by a private benefactor who
%                   has asked not to be named."
%   (c) omitted:   no compute sentence at all (and confirm the donor is content
%                   with that -- silence is also a choice they should get to make).
%
% Default if unreachable before the deadline: (b). It is honest, it discloses that
% the compute was donated rather than institutionally funded, and it exposes nobody.
%
% [GAP: compute acknowledgement -- BLOCKED on donor consent. See gate above.]

% [GAP: author list, affiliations and per-author contribution statement.
%  Deliberately NOT drafted yet -- see the authorship discussion, which is being
%  settled separately. Do not populate from git history alone.]

% Exactly ONE \bibliography command is allowed per document; the appendix must not
% carry a second one. \bibliographystyle is required -- the old AAAI class supplied it
% implicitly, stock article does not, which is why citations rendered as "??".
\bibliographystyle{plainnat}
\bibliography{refs}

\clearpage
\appendix

\section{Extraction Prompt and Schema}
\label{app:prompt}
% [GAP: paste PROMPT_EXTRACT from
%  intervention_graph_creation/src/prompt/final_primary_prompt.py, plus the node/edge
%  Pydantic schema summary. Trim to one page.]

\section{Deduplication Threshold Experiments}
\label{app:dedup}
% Reused verbatim from Gleb_overleaf_paper.txt Appendix D.

We evaluated four threshold configurations for node deduplication, measuring the
trade-off between precision (avoiding false merges) and recall (capturing
legitimate duplicates).

\begin{table}[h]
\centering
\small
\begin{tabular}{ccccc}
\toprule
\textbf{Cosine} & \textbf{Jaccard} & \textbf{Cand. Pairs} & \textbf{Dedup \%} & \textbf{Result} \\
\midrule
0.65 & 0.30 & 416 & 0.2\% & Too conservative \\
0.85 & 0.00 & 105,390 & 11.9\% & False clusters \\
\textbf{0.88} & \textbf{0.05} & \textbf{4,411} & \textbf{1.2\%} & \textbf{Selected} \\
0.95 & 0.00 & 503 & 0.2\% & Too conservative \\
\bottomrule
\end{tabular}
\caption{Deduplication threshold experiments. Candidate pairs shows node pairs
passing the similarity threshold; after Union-Find grouping fewer nodes are
removed as duplicates merge transitively (4{,}411 pairs yield 2{,}385 removed
nodes). The Jaccard-free configuration (0.85/0.00) created false clusters merging
unrelated concepts. \Cref{sec:claim3} reports the recall audit of the selected
configuration.}
\label{tab:dedup-thresholds}
\end{table}

\paragraph{Boundary examples (selected configuration).}
Representative merged pairs at the decision boundary (cosine $\approx$ 0.88):
\begin{itemize}
\setlength{\itemsep}{0pt}
\item ``US semiconductor export controls'' $\leftrightarrow$ ``AI chip export control regulations''
\item ``Misalignment of AI objectives'' $\leftrightarrow$ ``objective misalignment in smarter-than-human AI''
\item ``Insufficient AI safety research capacity'' $\leftrightarrow$ ``Insufficient AI alignment research capacity''
\end{itemize}

\paragraph{Edge type normalisation.}
The 12 canonical edge types after normalisation:
\textit{motivates, implemented\_by, validated\_by, caused\_by, mitigated\_by,
enabled\_by, addressed\_by, refined\_by, specified\_by, required\_by,
preceded\_by, related\_to}.

\section{Clustering Methods Comparison}
\label{app:clustering}
% Reused verbatim from Gleb_overleaf_paper.txt Appendix F.

\begin{table}[h]
\centering
\small
\begin{tabular}{lccc}
\toprule
\textbf{Method} & \textbf{Silhouette} & \textbf{CH Score} & \textbf{DB Index} \\
\midrule
Direct 1536D, $k=5$ & 0.022 & 3,831 & 5.18 \\
Direct 1536D, $k=8$ & 0.017 & 2,922 & 5.03 \\
UMAP 150D, $k=8$ & 0.287 & 97,380 & 1.27 \\
\textbf{UMAP 150D, $k=40$} & \textbf{0.298} & \textbf{55,410} & \textbf{1.15} \\
\bottomrule
\end{tabular}
\caption{Clustering quality metrics. Direct clustering in 1536D embedding space
fails (silhouette $\approx$ 0.02). UMAP reduction to 150D enables meaningful
topical clustering with a 13--17$\times$ silhouette improvement.
Calinski-Harabasz (CH) measures separation (higher is better); Davies-Bouldin
(DB) measures overlap (lower is better).}
\label{tab:clustering-methods}
\end{table}

\section{Topical Cluster Reference}
\label{app:clusters}
% Reused from Gleb_overleaf_paper.txt Appendix E. Reframed as a browsing index.
% [GAP: also publish cluster_representatives_20.json (20 representative nodes per
%  cluster) alongside the dataset so reviewers can audit cluster naming. Owner: Gleb.]

The 40 topical clusters from $k$-means on UMAP-reduced embeddings, cluster sizes
2{,}065--8{,}098 nodes. These are a \emph{browsing index} over the corpus, not a
taxonomy: co-membership indicates vocabulary proximity rather than a mechanistic
relationship, and boundaries between adjacent sub-fields overlap.

\begin{small}
\begin{enumerate}
\setlength{\itemsep}{0pt}
\item[0.] RL Impact Measures, Power \& Time-Bounded Optimization
\item[1.] Foundational Alignment Research \& AGI Safety
\item[2.] Deployment Monitoring \& Safety Evaluation
\item[3.] Training Optimization \& Hyperparameter Learning
\item[4.] Information Theory \& Communication Efficiency
\item[5.] Explainability \& Interpretability Methods
\item[6.] Compute Governance \& Hardware Regulation
\item[7.] AI Ethics \& Consciousness Research
\item[8.] Compute Scaling, Recursive Self-Improvement \& Brain Preservation
\item[9.] Value Learning \& Preference Aggregation
\item[10.] Fine-Tuning Methods \& Cognitive Empathy
\item[11.] Debate \& Iterated Amplification
\item[12.] AI Governance \& Legal Frameworks
\item[13.] Data Quality \& LLM Hallucinations
\item[14.] Reward Misspecification \& Proxy Gaming
\item[15.] Curriculum Learning \& Scientific Discovery Systems
\item[16.] Economic Impact \& Technological Unemployment
\item[17.] AI Safety Community \& Education
\item[18.] Transparency \& Accountability Mechanisms
\item[19.] Corrigibility \& Shutdown Safety
\item[20.] AI Forecasting \& Timeline Prediction
\item[21.] Wellbeing-Aware Recommenders, Misuse \& Content Moderation
\item[22.] Inner \& Outer Misalignment
\item[23.] Imitation Learning \& Sample Efficiency
\item[24.] Red Teaming \& Adversarial Testing
\item[25.] Adversarial Robustness \& OOD Generalization
\item[26.] Empirical Benchmarking \& Scaling Laws
\item[27.] Reward Modeling \& Human Feedback
\item[28.] Decision Theory \& Coherence
\item[29.] Activation Engineering \& Model Steering
\item[30.] Distribution Shift \& Robustness Failures
\item[31.] Formal Verification \& Proof Systems
\item[32.] Containment \& Sandboxing Strategies
\item[33.] Biosecurity \& Critical Infrastructure
\item[34.] Legible Motion Planning, Reachability \& Safe Delegation
\item[35.] Deceptive Alignment \& Treacherous Turn
\item[36.] AI Risk Perception, Governance \& Research Prioritization
\item[37.] Funding \& Resource Allocation
\item[38.] Existential Risk \& Catastrophic Scenarios
\item[39.] Public Outreach, Framing \& Community Building
\end{enumerate}
\end{small}

\section{Similarity Threshold and Hop Sensitivity}
\label{app:sensitivity}
% Reused from Gleb_overleaf_paper.txt Appendices G + H, merged.

\paragraph{Similarity threshold ($\tau$).}
We tested $\tau \in \{0.80, 0.85, 0.90, 0.95\}$ on within-category node pairs.
At $\tau{=}0.80$, 169{,}083 similarity edges connect conceptually related nodes
across documents. Higher thresholds progressively fragment the graph:
$\tau{=}0.85$ reduces edges by approximately 50\%, $\tau{=}0.90$ by over 90\%,
and $\tau{=}0.95$ captures only near-duplicate phrasings. The main text uses
EDGE-only paths for all quantitative claims precisely because of this
sensitivity.

\paragraph{Similarity hops.}
Graph traversal with 0--3 similarity hops on 500 sampled risks:

\begin{table}[h]
\centering
\footnotesize
\begin{tabular}{cccc}
\toprule
\textbf{Hops} & \textbf{Scope (mean)} & \textbf{Intv. (mean)} & \textbf{Intv. (med.)} \\
\midrule
0 & 14 & 2.5 & 2 \\
1 & 300 & 56 & 10 \\
2 & 6,209 & 1,172 & 120 \\
3 & 26,827 & 5,065 & 2,384 \\
\bottomrule
\end{tabular}
\caption{Similarity hop comparison. Scope is the mean number of nodes reachable
from a risk via BFS; Intv. is reachable interventions. Each hop multiplies reach
by roughly 20$\times$, which is why hop-augmented reachability appears here as a
sensitivity rather than as evidence in the main text.}
\label{tab:sim-hops}
\end{table}

\paragraph{PageRank mixing coefficient ($\alpha$).}
For the centrality analyses in \cref{sec:claim3} we combine structural and
similarity transition matrices as
$P = \alpha P_{\text{struct}} + (1-\alpha) P_{\text{sim}}$.
Rankings are stable across $\alpha \in [0.5, 0.9]$ (Spearman $\rho = 1.0$,
identical top-20); only pure structural weighting ($\alpha = 1.0$) deviates, with
2 risks leaving the top-20. This stability is a robustness property of the mixing
parameter and says nothing about the merge artefact documented in
\cref{sec:claim3}: the rankings are stably contaminated.

\section{Race-Framing Classifier Validation}
\label{app:race}
% Reused from Gleb_overleaf_paper.txt Appendix J, kept because it is a valid
% classifier validation; prevalence claims moved to the main text.

We validated keyword-based race-framing classification (keywords: ``competi-'',
``race'') on a stratified sample of 52 problem-analysis nodes.
\textbf{Stratum 1} ($n{=}36$): all race-flagged PA nodes from top-100 single-path
risks --- 36/36 = 100\% precision. Example PAs: ``Competitive deployment
pressures reducing safety measures,'' ``AI arms race dynamics reducing safety
investment.'' \textbf{Stratum 2} ($n{=}16$): random sample from remaining risks
--- 12/16 = 75\% precision. The four false positives matched keywords but
described competition \emph{among} AI systems rather than a human race \emph{to
develop} AI. Overall precision 48/52 = 92\%.

\paragraph{Recall estimation.} On a stratified sample ($n{=}50$) of non-flagged
PA nodes, one false negative appeared among top-100 single-path risks (``AI
developers racing without adequate safety''), giving recall 36/37 = 97.3\% in
that stratum. Recall falls to an estimated $\sim$68\% in tiers 101--500 due to
morphological variants and synonyms without keywords.

\paragraph{Scope.} This validates the \emph{classifier} on the sampled strata and
establishes nothing about prevalence. \Cref{sec:claim3} reports corpus-wide
prevalence (2.21\% of problem-analysis nodes), measured separately by exhaustive
scan.

\section{A Failure Case}
\label{app:failure}

Quality cuts do not make every extracted chain meaningful. The chain below passes
every structural filter we apply --- seven nodes, all five reasoning stages, edge
confidence $\geq 3$, a maturity-4 intervention endpoint, single source --- and is
nonetheless spurious. Its source is an expository mathematics article proving
that there are infinitely many primes.

\begin{quote}\small
\textbf{risk} Prime scarcity affecting cryptographic key space $\rightarrow$
\textbf{pa} Finite prime assumption in key generation algorithms $\rightarrow$
\textbf{ti} Contradiction from finite prime list assumption $\rightarrow$
\textbf{dr} Product-plus-one construction to find new prime $\rightarrow$
\textbf{im} Compute $P = \prod(\text{all known primes}) + 1$ $\rightarrow$
\textbf{va} Proof that $P$ has a prime divisor outside the original list
$\rightarrow$ \textbf{intv} Integrate infinite-prime theorem into cryptographic
key generation standards
\end{quote}
% SRC: node ids [60117..60123], source https://arbital.com/p/infinitely_many_primes

The extraction has taken Euclid's proof and cast it into the risk--intervention
schema, inventing a safety framing (``prime scarcity'') the source never asserts.
This is the characteristic failure mode of schema-guided extraction over a
heterogeneous corpus: when a document contains a problem, an argument and a
construction, the schema will find a chain whether or not the document concerns
safety at all. Note which failure this is: the source document is a real
mathematical argument, so the failure is not that the pipeline read an
out-of-domain document, but that it \emph{invented} a risk framing absent from
the source. That is a fidelity failure of the kind \cref{sec:claim1} identifies
as needing a control, and it is why the structural statistics there are reported
as structural properties, never as evidence that a chain is substantively
correct. Establishing the rate at which this happens requires the schema-ablation
and degraded-source controls named in \cref{sec:claim1}; this appendix is an
existence proof only.

% [GAP: optionally add 1-2 further chains (e.g. one with a missing stage) if space
%  permits. The failure case above is the one that must stay.]

% NOTE: no \bibliography here -- it lives once, before the appendix. A second one
% makes BibTeX abort with "Illegal, another \bibdata command".

\end{document}
